\chapter{Conclusioni}

Il progetto S.I.D.E. ha dimostrato l’efficacia di un ambiente simulato per sistemi ICS/SCADA, capace di generare traffico realistico e testare la resilienza dei protocolli Modbus TCP, Siemens S7 e OPC-UA. L’implementazione di server e client simulati ha permesso di osservare in tempo reale il comportamento dei sistemi e degli IDS, con gestione di anomalie controllate e monitoraggio accurato.

\paragraph{Sviluppi futuri}
Tra le principali direzioni di evoluzione del progetto si segnalano:
\begin{itemize}
\item \textbf{Automazione CI/CD}: creazione di ambienti simulati containerizzati (Docker) con configurazioni di rete variabili per testare singoli protocolli per volta in vari scenari SCADA in pipeline automatizzate;
\item \textbf{Attacchi controllati dal front-end}: possibilità di generare scenari offensivi direttamente dall’interfaccia utente per valutare l’efficacia delle difese e il comportamento degli IDS;
\item \textbf{Machine Learning avanzato}: integrazione di modelli sofisticati e dataset estesi sui protocolli industriali per migliorare la rilevazione delle anomalie e la previsione di attacchi complessi. Ad esempio : GNN ( Graph Neural Network ) 
\end{itemize}